\documentclass[aps,prb,twocolumn,10pt,nofootinbib,superscriptaddress]{revtex4-2}
\usepackage{amsmath,amssymb,amsthm,mathtools,bm}
\usepackage{booktabs}
\usepackage[colorlinks=true,linkcolor=blue,citecolor=blue,urlcolor=blue]{hyperref}
\usepackage{microtype}
\hypersetup{pdftitle={Exact many-body relaxation gap of a fermionic reset circuit},
  pdfauthor={Class-A fermionic adaptive-circuit project}}
\newcommand{\Id}{\hat{\mathbf{1}}}
\newcommand{\id}{\mathbf{1}}
\newcommand{\Tr}{\operatorname{Tr}}
\newcommand{\spec}{\operatorname{spec}}
\newcommand{\E}{\widetilde{\mathcal E}}

\begin{document}
\title{Exact many-body relaxation gap of a fermionic reset circuit}
\author{Class-A fermionic adaptive-circuit project}
\noaffiliation
\date{October 1, 2026}
\begin{abstract}
For a fixed sequence of perfect occupation resets, we prove that the
charge-neutral many-body channel gap equals the covariance decay rate
$-2\log\rho(A)$, where $A$ is the ordered product of single-particle
complement projectors. The proof uses an exact triangular hierarchy of
fermionic operator degrees and requires neither Gaussian input states nor
commuting measured modes. We distinguish this physical relaxation sector
from the unrestricted operator-space gap.
\end{abstract}
\maketitle
\setcounter{tocdepth}{1}
\tableofcontents

\section{Reset channel and neutral sector}
Consider $N$ complex modes with
$\{\hat\psi_i,\hat\psi_j^\dagger\}=\delta_{ij}$ and
$\{\hat\psi_i,\hat\psi_j\}=0$. At step $j$, measure the normalized mode
$\hat\chi_j=\sum_i v_{ji}\hat\psi_i$, $v_j^\dagger v_j=1$, and reset its
occupation $\hat N_j=\hat\chi_j^\dagger\hat\chi_j$ to a prescribed
$s_j\in\{0,1\}$. After discarding the outcome and ancillary mode, the
specified Kraus channel is
\begin{align}
 \E_{j,0}(\hat\rho)&=\hat\chi_j\hat\rho\hat\chi_j^\dagger
 +(\Id-\hat N_j)\hat\rho(\Id-\hat N_j),\nonumber\\
 \E_{j,1}(\hat\rho)&=\hat\chi_j^\dagger\hat\rho\hat\chi_j
 +\hat N_j\hat\rho\hat N_j. \label{eq:reset}
\end{align}
The Kraus completeness relation follows from $\hat N_j^2=\hat N_j$.
With $\mathcal D[\hat L](\hat\rho)=\hat L\hat\rho\hat L^\dagger
-\{\hat L^\dagger\hat L,\hat\rho\}/2$, these are exactly
$\mathrm{Id}+\mathcal D[\hat\chi_j]+\mathcal D[\hat N_j]$ and
$\mathrm{Id}+\mathcal D[\hat\chi_j^\dagger]+\mathcal D[\hat N_j]$,
respectively, rather than continuous-time exponentials.

One cycle is the composition
$\E=\E_{M,s_M}\circ\cdots\circ\E_{1,s_1}$, with step 1 acting first;
the sequence repeats every cycle. Set
\begin{equation}
 P_j=v_jv_j^\dagger,\qquad Q_j=\id-P_j,\qquad
 A=Q_M\cdots Q_1. \label{eq:A}
\end{equation}
The modes may be nonorthogonal. In particular, overcomplete Wannier (OW)
stabilizers need not commute~\cite{BPJ}; their order is retained in $A$.

For $\overline G_{ab}=\Tr(\overline{\hat\rho}\,
\hat\psi_a^\dagger\hat\psi_b)$, direct application of
Eq.~\eqref{eq:reset} gives
\begin{align}
 \overline G'&=Q_j\overline GQ_j+s_jP_j,\nonumber\\
 \overline G_{n+1}&=A\overline G_nA^\dagger+B,
 \qquad \delta G_{n+1}=A\delta G_nA^\dagger. \label{eq:cov}
\end{align}
Here $B$ is independent of the input and
$\delta G=\overline G-G_{\rm ss}$. Thus, for eigenvalues $a_i$ of $A$,
the covariance multipliers are $a_i a_j^*$. Define
\begin{equation}
 \begin{gathered}
 r=\rho(A)=\max_i|a_i|,\qquad 0<r<1,\\
 \Delta_C=-2\log r.
 \end{gathered} \label{eq:rate}
\end{equation}
All rates below are per complete cycle.

The relevant operator sector is
\begin{equation}
 \mathcal V_0=\{\hat O:[\hat N_F,\hat O]=0\},\qquad
 \hat N_F=\sum_i\hat\psi_i^\dagger\hat\psi_i. \label{eq:neutral}
\end{equation}
Each Kraus operator has definite charge, so $\E$ is $U(1)$ covariant
and preserves $\mathcal V_0$. Density matrices of fixed-number Slater
determinants and mixtures of number sectors belong to $\mathcal V_0$.
Fill/empty events change particle number while preserving this property.
Neutrality is an operator-space condition, not a fixed-filling constraint.

\section{Exact spectral factorization}
\textit{Local deletion rule.} Choose an orthonormal mode basis beginning
with $\hat\chi_j$ and let $\hat B$ involve only the other modes. From
Eq.~\eqref{eq:reset}, for an even total operator,
\begin{align}
 \E_{j,s}^\dagger(\hat B)&=\hat B &&(\hat B\text{ even}),\nonumber\\
 \E_{j,s}^\dagger(\hat\chi_j^{(\dagger)}\hat B)&=0
 &&(\hat B\text{ odd}),\nonumber\\
 \E_{j,s}^\dagger(\hat N_j\hat B)&=s\hat B
 &&(\hat B\text{ even}). \label{eq:rule}
\end{align}
For example, filling gives
$\hat\chi_j\hat N_j\hat B\hat\chi_j^\dagger
+\hat N_j\hat N_j\hat B\hat N_j
=(\Id-\hat N_j)\hat B+\hat N_j\hat B=\hat B$;
both terms vanish for depletion. A single measured creation or
annihilation operator has zero Kraus sandwiches.

\textit{Invariant filtration.} A basis of $\mathcal V_0$ consists of
normal-ordered monomials
\begin{equation}
 \hat O_{I,J}=\hat\psi_{i_1}^\dagger\cdots\hat\psi_{i_p}^\dagger
 \hat\psi_{j_p}\cdots\hat\psi_{j_1},\qquad |I|=|J|=p,
 \label{eq:monomials}
\end{equation}
including $p=0$. Equation~\eqref{eq:rule} either preserves $(p,p)$,
annihilates the monomial, or lowers it to $(p-1,p-1)$. A change of
one-particle basis preserves these degrees. Hence the span of all
monomials through degree $(p,p)$ is invariant under $\E^\dagger$.

On the quotient by lower degrees, each index is projected by $Q_j$ or
$Q_j^*$. Explicitly, for the moments
$F_{I,J}=\Tr(\hat\rho\hat O_{I,J})$, the degree-preserving block of one
reset has entries
\begin{equation}
 H_{j,p}[I,J;K,L]
 =\det Q_j[I,K]\,\det Q_j[J,L]^*. \label{eq:minors}
\end{equation}
Here $Q_j[I,K]$ is a $p\times p$ submatrix. These determinants are the
matrix entries of exterior powers. The Cauchy--Binet identity gives
$\wedge^p(XY)=(\wedge^p X)(\wedge^p Y)$; applying it in the actual reset
order yields the homogeneous moment propagator over one cycle,
\begin{equation}
 T_p=(\wedge^p A)\otimes(\wedge^p A^*),\qquad
 p=0,\ldots,N. \label{eq:blocks}
\end{equation}
The dual operator-coefficient block is its transpose, with the same
spectrum. Reset-source terms couple only to lower degrees and therefore
do not change the eigenvalues of this block-triangular hierarchy.

The eigenvalues of an exterior power are products of distinct
one-particle eigenvalues. Consequently, counting algebraic multiplicity,
\begin{equation}
 \spec(\E|_{\mathcal V_0})=
 \bigcup_{p=0}^{N}
 \left\{\prod_{i\in I}a_i\prod_{j\in J}a_j^*:
 |I|=|J|=p\right\}. \label{eq:spectrum}
\end{equation}
The sets $I,J$ are independent and may overlap; each index occurs at
most once within either set. This accounts for
$\sum_p\binom Np^2=\binom{2N}{N}=\dim\mathcal V_0$ eigenvalues.
Schur triangularization proves the exterior-power statement even when
$A$ has Jordan blocks. No Gaussian-state assumption is used; Gaussian
maps provide a related general framework~\cite{Bravyi}.

\textit{Gap.} The $p=0$ eigenvalue is 1. Every $p\geq1$ product satisfies
$|\Lambda|\leq r^{2p}\leq r^2$. Equality is attained at $p=1$ by
$I=J=\{i_*\}$ with $|a_{i_*}|=r$. The stationary eigenvalue is therefore
simple in this sector and
\begin{equation}
 \boxed{\Delta_{q=0}\equiv-\log
 \max_{\Lambda\ne1}|\Lambda|=-2\log r=\Delta_C.}
 \label{eq:gap}
\end{equation}
This proves equality, strengthening the upper bound obtained from
quadratic spectral inclusion alone.

\section{Parity sectors and scope}
Allow $p$ creations and $\ell$ annihilations in Eq.~\eqref{eq:monomials},
with net operator charge $q=p-\ell$.
For $p+\ell$ even, the same proof gives blocks
$(\wedge^p A)\otimes(\wedge^\ell A^*)$. Every nonconstant product contains
at least two factors, and the neutral block attains $r^2$. Thus
$\Delta_{\rm even}=\Delta_{q=0}$; even means commuting with fermion
parity, not occupying only one Hilbert-space parity block.

For completeness, the unrestricted spectrum of the specified Kraus
channel requires an additional step. Let $\hat\Pi=(-1)^{\hat N_F}$ and
$\mathcal S(\hat O)=\hat\Pi\hat O$, so $\mathcal S^2=\mathrm{Id}$.
For even and odd Kraus operators $\hat K_{\rm e},\hat K_{\rm o}$,
\begin{equation}
 (\mathcal S\E_j^\dagger\mathcal S)(\hat O)
 =\hat K_{\rm e}^\dagger\hat O\hat K_{\rm e}
 -\hat K_{\rm o}^\dagger\hat O\hat K_{\rm o}. \label{eq:twist}
\end{equation}
On odd monomials, this sign restores Eq.~\eqref{eq:rule} with the
parities of $\hat B$ reversed: moving an odd $\hat B$ through the odd
Kraus operator cancels the minus sign. Similarity preserves spectra and
composition, giving the same blocks with $p+\ell$ odd. Their maximal
modulus is $r$, attained at total degree one. Hence
\begin{equation}
 \Delta_{\rm full}=-\log r=\tfrac12\Delta_C,
 \qquad \Delta_{\rm even}=\Delta_{q=0}=\Delta_C. \label{eq:sectors}
\end{equation}
Physical density matrices have even parity, and the preparations under
discussion are also charge neutral. Their channel relaxation is governed
by Eq.~\eqref{eq:gap}. The proof concerns a fixed finite system, exact
resets, and a repeated deterministic schedule. It establishes an
asymptotic spectral rate, allowing Jordan prefactors; it does not by itself
bound a size-uniform mixing time or a conditional-trajectory purification
time. Standalone number dephasing, imperfect feedback, and averaging
different schedules require a new higher-order analysis.

\begin{thebibliography}{2}
\bibitem{BPJ}
A. Bhuiyan, H. Pan, and C.-M. Jian,
\href{https://doi.org/10.1103/nk38-ygyq}{Phys. Rev. Research \textbf{8}, 023147 (2026)}.
\bibitem{Bravyi}
S. Bravyi, Lagrangian representation for fermionic linear optics,
\href{https://arxiv.org/abs/quant-ph/0404180}{Quantum Inf. Comput. \textbf{5}, 216 (2005)}.
\end{thebibliography}
\end{document}
